\documentclass[11pt,a4paper]{article}
\usepackage{amsmath,amssymb,amsthm}
\usepackage[margin=1in]{geometry}

\newtheorem{theorem}{Theorem}
\newtheorem{lemma}{Lemma}
\newtheorem{corollary}{Corollary}
\newtheorem{proposition}{Proposition}

\DeclareMathOperator{\re}{Re}
\DeclareMathOperator{\im}{Im}

\title{The First-Order $M$-Identity}
\author{Working notes}
\date{March 2026}

\begin{document}
\maketitle

\section{The identity}

\begin{theorem}[First-order $M$-identity]\label{thm:identity}
Define $g'(t) = i\,\dfrac{\xi'}{\xi}(\tfrac{1}{2}+it)
+ t\,\left(\dfrac{\xi'}{\xi}\right)'(\tfrac{1}{2}+it)$.
Then:
\begin{equation}\label{eq:identity}
  g'(t) = \sum_\rho \frac{4t^3}{(t^2 + (\rho-\tfrac{1}{2})^2)^2},
\end{equation}
where the sum runs over non-trivial zeros $\rho$ of $\zeta$,
paired via the functional equation (each pair $\rho, 1-\rho$
contributes one term with
$w_\rho^2 = (\rho - \tfrac{1}{2})^2$).
\end{theorem}

\begin{proof}
From the Hadamard product
$\xi(s) = \xi(\tfrac{1}{2})\prod_\rho
(1 - (s-\tfrac{1}{2})^2/(\rho-\tfrac{1}{2})^2)$:
\[
  \frac{\xi'}{\xi}(s)
  = \sum_\rho \frac{2(s-\tfrac{1}{2})}{(s-\tfrac{1}{2})^2
    - (\rho-\tfrac{1}{2})^2}.
\]
Setting $u = s - \tfrac{1}{2}$ and $w^2 = (\rho-\tfrac{1}{2})^2$:
\begin{align*}
  \frac{\xi'}{\xi} &= \sum \frac{2u}{u^2 - w^2}, \\[4pt]
  \left(\frac{\xi'}{\xi}\right)'
  &= \sum \frac{-2(u^2 + w^2)}{(u^2 - w^2)^2}.
\end{align*}
At $s = \tfrac{1}{2} + it$, $u = it$:
\[
  i \cdot \frac{2it}{-t^2 - w^2}
  + t \cdot \frac{-2(-t^2 + w^2)}{(-t^2-w^2)^2}
  = \frac{-2t}{t^2+w^2}
  + \frac{2t(t^2-w^2)}{(t^2+w^2)^2}
  = \frac{4t^3}{(t^2+w^2)^2}. \qedhere
\]
\end{proof}

\section{Consequences}

\begin{corollary}[On-line positivity]
If all non-trivial zeros satisfy $\re(\rho) = 1/2$ (RH), then
$(\rho-1/2)^2 = -\gamma_\rho^2$ and
\begin{equation}
  g'(t) = \sum_{\gamma > 0}
    \frac{4t^3}{(t^2-\gamma^2)^2}
  = \sum_{\gamma > 0}
    \frac{\gamma}{(t-\gamma)^2} \cdot \frac{4t^3}{(t+\gamma)^2}
  > 0.
\end{equation}
Each term is positive. The positivity is manifest.
\end{corollary}

\begin{corollary}[Connection to $M$]
$M(\tfrac{1}{2}+\epsilon+it) = \epsilon\,g'(t) + O(\epsilon^2)$.
Therefore $g'(t) > 0$ is a \emph{necessary condition} for
$M(s) > 0$ near the critical line, hence necessary for RH.
\end{corollary}

\begin{proposition}[Off-line contribution]
For a zero $\rho = \beta + i\gamma$ with
$\alpha = \beta - 1/2 \neq 0$, the contribution to the real
part of $g'(t)$ from the pair $(\rho, 1-\rho)$ is
\begin{equation}
  \re\!\left[\frac{4t^3}{(t^2+w^2)^2}\right]
  = \frac{4t^3(A^2-B^2)}{(A^2+B^2)^2},
\end{equation}
where $A = t^2+\alpha^2-\gamma^2$ and $B = 2\alpha\gamma$.
This is negative when $B^2 > A^2$, i.e., when
$4\alpha^2\gamma^2 > (t^2+\alpha^2-\gamma^2)^2$.
\end{proposition}

\section{The open question}

$g'(t) > 0$ is equivalent to: the on-line zero contributions
(always positive) outweigh any off-line contributions
(possibly negative).

Conrey (1989) proved $\geq 40\%$ of zeros are on the
critical line. The question: does this $40\%$ guarantee
dominate? The on-line terms scale as $(t-\gamma)^{-2}$
(strong near each zero), while off-line terms have the
same scaling but with a negative sign for certain
$(t,\gamma,\alpha)$ configurations.

A proof that $g'(t) > 0$ unconditionally would establish a
necessary condition for RH using only zero density estimates
and the Hadamard product identity.

\end{document}
